\documentclass[11pt]{article}
\usepackage[margin=1in]{geometry}
\usepackage{amsmath,amssymb,amsthm,mathtools}

\newtheorem{theorem}{Theorem}[section]
\newtheorem{lemma}[theorem]{Lemma}
\newtheorem{proposition}[theorem]{Proposition}
\newtheorem{corollary}[theorem]{Corollary}
\theoremstyle{definition}
\newtheorem{definition}[theorem]{Definition}
\theoremstyle{remark}
\newtheorem{remark}[theorem]{Remark}
\newcommand{\J}{\mathcal J}
\newcommand{\cG}{\mathcal G}
\newcommand{\cL}{\mathcal L}
\newcommand{\C}{\mathbb C}
\newcommand{\R}{\mathbb R}
\newcommand{\N}{\mathbb N}
\newcommand{\dd}{\,\mathrm d}

\title{A Two-Thirds Hyperbolicity Wedge for Jensen Polynomials of
Riemann's Xi-Function\\\large Reader's synopsis and proof map}
\author{John Savva\\{\small Independent researcher}}
\date{Version 1.0 \textemdash{} 21 August 2026}

\begin{document}
\maketitle

\begin{center}
\fbox{\parbox{0.91\textwidth}{\small
\textbf{Purpose and review status.} This is a compact synopsis; the maintained
scholarly sources are the main paper and technical supplement in the
\texttt{paper} directory.  The user-executed exact Mathematica reconstruction
completed the independent-CAS checks M1--M4.  Lean now checks the complete
T1--T5 analytic chain, the exact xi branch through six matches, and the
critical-radius conclusion under narrow typed literature inputs.  The chosen
terminal endpoint also constructs the final xi-specific sixth-multiplier
certificate, connects it to the actual transformed Jensen polynomial, and
checks a single global exact-root conclusion conditional only on those typed
inputs.  The MSS interface requires positive lower endpoints, and Lean derives
them from the displayed $256d$ geometry before applying the imported result.
All available reviews are AI-assisted evidence, not human or peer review.}}
\end{center}

\begin{abstract}
Let $\gamma(n)$ be the Taylor coefficients of Riemann's completed
$\xi$-function in the standard even normalization, and let
\[
  \J^{d,n}(X)=\sum_{j=0}^d \binom{d}{j}\gamma(n+j)X^j.
\]
We prove that there is an absolute $K>0$ such that
$n^2\log(n+2)\ge Kd^3$ implies that $\J^{d,n}$ has $d$ distinct negative
real zeros.  The gain from the earlier three-fifths architecture comes from
matching six logarithmic coefficients with a four-parameter product of two
Jacobi factors.  The proof combines a direct sectorial saddle analysis for
the xi coefficients, a sixth-order residual estimate, finite-free
convolution, and a non-perturbative critical-point recurrence.  We state the
trust boundary precisely: decisive symbolic identities have exact
SymPy/Lean regression evidence and an exact user-executed Mathematica
second-CAS reconstruction, and the complete T1--T5 analytic chain is
kernel-checked in Lean.  The xi-specific multiplier, its six nodes and unit
bound, interval certificate, and headline assembly are also kernel checked.
The remaining formal boundary consists only of typed Jacobi, MMP, and MSS
literature inputs.
This synopsis compresses that proof; detailed inequalities and evidence maps
are in the main paper and supplement.
\end{abstract}

\section{Statement and architecture}

\begin{theorem}[Two-thirds wedge]\label{thm:main}
There is an absolute constant $K>0$ such that for integers $d\ge1$ and
$n\ge0$, the implication
\[
             n^2\log(n+2)\ge Kd^3
\]
ensures that $\J^{d,n}$ has exactly $d$ distinct negative real zeros.
\end{theorem}

The proof is developed block by block below and is closed explicitly in
Section~\ref{sec:assembly}.  The synopsis uses ordinary section-based theorem
numbers; the labels T1--T18 in the expanded paper and assurance matrix are
stable evidence identifiers rather than additional theorem numbers.  Thus
the sectorial coefficient theorem is T5 and the present main theorem is T18.

The proof is effective in principle after tracing the saddle and branch
constants, but the present unweighted sup-norm contraction yields an
astronomically large sufficient threshold.  We make no claim that the
resulting constant is computationally useful, and we do not print a
numerical value of $K$.

The proof has four blocks.  First, Sections~\ref{sec:mellin}--\ref{sec:deriv}
derive a sectorial coefficient asymptotic and its logarithmic derivatives
through order six.  Second, Sections~\ref{sec:model}--\ref{sec:branch} choose
a positive four-parameter comparison model matching six coefficients.
Third, Sections~\ref{sec:finitefree}--\ref{sec:recurrence} establish positive,
simple, localized roots and the derivative radius of that comparison.
Finally, Sections~\ref{sec:residual}--\ref{sec:assembly} transfer signs to the
actual Jensen polynomial.

Throughout, a statement called ``uniform'' is uniform after the fixed
nested sectors and compact parameter box have been chosen.  Holland's main
theorem is not a premise.  His paper supplies architecture and several
formulas which are rederived or connected directly to their primary
sources.

\section{The Mellin coefficient identity}\label{sec:mellin}

Put
\[
 \Theta(t)=\sum_{m\ge1}e^{-\pi m^2t},\qquad
 \omega(t)=\tfrac12\bigl(2t^2\Theta''(t)+3t\Theta'(t)\bigr).
\]
Riemann's integral and the substitution $t=e^{2u}$ give an even Fourier--
Laplace representation of $\Xi$.  Differentiating at the origin, retaining
both halves of the even integral, gives the exact normalization
\begin{equation}\label{eq:factor8}
  \gamma(n)=\frac{8\,n!}{(2n)!}
       \int_0^\infty \omega(e^{2u})e^{u/2}u^{2n}\,\dd u .
\end{equation}
The factor eight in \eqref{eq:factor8} is important when the main term is
compared character for character with the printed GORZ/GORTTW convention.
It cancels from logarithmic derivatives but not from a source-fidelity
check.  Indeed Riemann's kernel in these variables is
$8e^{u/2}\omega(e^{2u})$, and extracting the coefficient of $w^{2n}$
from $\cosh(wu)$ contributes $1/(2n)!$.  Since our convention is
$\xi(1/2+w)=\sum_{n\ge0}\gamma(n)w^{2n}/n!$, this gives
\eqref{eq:factor8}.  Equivalently, the prefactor is
$8\sqrt\pi/(4^n\Gamma(n+1/2))$ by duplication.  No asymptotic input enters.

Expanding $\Theta$ separates the $m=1$ mode from $m\ge2$.  If $M$ denotes the
complex Mellin exponent, the leading phase has one saddle and the remaining
theta modes acquire an additional negative exponential proportional to
$m^2-1$.  This exact decomposition is the starting point for the sectorial
analysis.

\section{The sectorial saddle variable}\label{sec:saddle}

For $0<\theta<\pi/2$, use principal logarithms in
$|\arg N|\le\theta$, and set
\[
 \Psi(L)=L\left(\pi e^L+\frac34\right),\quad
 w=\log N,\quad L_0=w-\log w-\log\pi.
\]

\begin{lemma}[Sectorial saddle, Lemma S]\label{lem:saddle}
For $|N|\ge e^{12}$ and $|\arg N|\le\theta$ there is a unique holomorphic
branch $L_N$, agreeing with the positive real solution, such that
$N=\Psi(L_N)$ and $|L_N-L_0|<1$.  Uniformly on every fixed closed smaller
sector,
\[
 |L_N|\asymp\log|N|,\qquad r:=L_N^{-1}\longrightarrow0,\qquad
 \sigma:=L_N/N\longrightarrow0.
\]
Moreover $|L_N|\ge\tfrac12\log|N|$, $|r|,|\sigma|\le7/50$, and
\begin{align}
 Q_N&=(1+L_N)N-\tfrac34L_N^2
     =NL_N(1+r-\tfrac34\sigma),\label{eq:Q}\\
 |1+r-\tfrac34\sigma|&\ge9/16,\qquad L_N'=L_N/Q_N.\label{eq:Qbound}
\end{align}
\end{lemma}

\begin{proof}
Consider
\[
 G(L)=L+\log L+\log\pi-
       \log\left(1-\frac{3L}{4N}\right)-\log N
\]
on $|L-L_0|\le1$.  If $\ell=\log|N|\ge12$, define
\[
 m(\ell)=\ell-\log(\ell+1)-\log\pi-1,
 \qquad
 M_\theta(\ell)=\ell+2+\log(\ell+1)+\theta/\ell+\log\pi.
\]
Then
\[
 \Re L\ge m(\ell)>0
\]
and
\[
 |L|\le M_\theta(\ell).
\]
Thus all principal logarithms are legal and
$|3L/(4N)|<10^{-4}$.  On the boundary, compare $G$ with
$H(L)=L-L_0$.  The elementary estimate $|\log(1+z)|\le2|z|$ for
$|z|\le1/2$ gives
\[
 \sup_{\partial D}|G-H|
 \le 2\frac{\log(\ell+1)+\theta/\ell+\log\pi+1}{\ell}
   +\frac{3M_\theta(\ell)}{2e^\ell}<1.
\]
Rouche's theorem gives one simple zero.  Exponentiating at that zero yields
$N=\Psi(L)$.  Simplicity and the holomorphic implicit-function theorem patch
the local zeros into a branch; conjugation and real uniqueness identify the
positive real branch.  The displayed disc bounds give the asserted
asymptotics.  In particular $m(\ell)>\ell/2$ for $\ell\ge12$, so
$|L_N|\ge\Re L_N\ge\ell/2$.

For \eqref{eq:Qbound}, the reverse triangle inequality and
$|r|,|\sigma|\le7/50$ give
\[
 |1+r-\tfrac34\sigma|
 \ge1-\frac7{50}-\frac34\frac7{50}
 =\frac{151}{200}>\frac9{16}.
\]
Hence $Q_N\ne0$ before implicit differentiation is invoked.  Differentiating
$N=\Psi(L_N)$ then gives $L_N'=L_N/Q_N$.
\end{proof}

The same exact argument gives
\begin{equation}\label{eq:reducedden}
 |4+4r-3\sigma|\ge4-4\frac7{50}-3\frac7{50}=\frac{151}{50}.
\end{equation}
This is the reduced denominator margin used at derivative order six.

\section{Contour localization and the coefficient theorem}\label{sec:coefficient}

\begin{theorem}[Sectorial coefficient asymptotic]\label{thm:sector}
Fix $\theta_0=1/400$ and $\theta_1=1/200$.  On the inner sector, the analytic
continuation of the xi coefficient moment equals its explicit saddle main
term times
\[
                  1+O\left(\frac{\log|M|}{|M|}\right),
\]
uniformly as $|M|\to\infty$.  The main term is the one obtained from
\eqref{eq:factor8} and agrees, after the factor-eight normalization, with the
printed GORZ main term.  It is nonzero on the inner sector.
\end{theorem}

\begin{proof}
The leading $m=1$ Mellin integral is moved, in $u=\log t$, to the horizontal
ray through the saddle supplied by Lemma~\ref{lem:saddle}.  The move stays in
the principal-log domain.  For $|\arg s|\le1/100$, the logarithmic saddle
equation and $\Re L_s\ge m(12)>7.29$ give $|\Im L_s|<1/20$.  Writing
\[
 K_s=\frac{s}{L_s}\left(1+\frac1{L_s}-\frac{3L_s}{4s}\right),
\]
the corresponding argument estimate gives
$|\arg K_s|<1/20$ and
$\Re K_s\ge\cos(1/20)|K_s|$.  This right-half-plane bound---not merely the
nonvanishing consequence of \eqref{eq:Qbound}---controls the Gaussian and
the strict concavity on the shifted ray.  The two endpoint connectors are
$\mathcal A(s)O(e^{-c|s|\log\log|s|})$, where $\mathcal A(s)$ is the saddle
main term.  Splitting into a central Gaussian window and two tails gives a
relative error $O(|K_s|^{-1})$; the signed cubic Gaussian moment removes the
otherwise larger absolute error.

For $m\ge2$, the same contour is legal and the extra theta exponent is
uniformly negative.  Summing the modes is therefore exponentially smaller
than the first mode.  Sectorial Stirling, applied on fixed nested sectors,
then yields the asserted main term and relative error.  The detailed
connector, central-window, signed-cubic, tail, mode-summation, and Stirling
estimates are expanded in the main paper and technical supplement.

The proof just described is carried out with fixed sector/radius constants;
there is no assertion that every later error constant is uniform for an
arbitrary free $\theta\in(0,\pi/2)$.
\end{proof}

\section{Logarithmic derivatives through order six}\label{sec:deriv}

Let $h(x)$ be the logarithm of the exact moment normalization on the positive
axis and let $\cL_x=L_{2x-2}$.  Cauchy's inequality on proportional discs
inside the two fixed sectors differentiates the relative error in
Theorem~\ref{thm:sector}.  Total implicit differentiation of the explicit
main term gives
\begin{align}
 h''(x)&=\frac{2}{x\cL_x}(1+O(\cL_x^{-1})),\nonumber\\
 h^{(3)}(x)&=-\frac{2}{x^2\cL_x}(1+O(\cL_x^{-1})),\nonumber\\
 h^{(4)}(x)&=\frac{4}{x^3\cL_x}(1+O(\cL_x^{-1})),\label{eq:tower}\\
 h^{(5)}(x)&=-\frac{12}{x^4\cL_x}(1+O(\cL_x^{-1})),\nonumber\\
 h^{(6)}(x)&=\frac{48}{x^5\cL_x}(1+O(\cL_x^{-1})).\nonumber
\end{align}
Only the bound in the last line is needed.  Before the chain rule $N=2x-2$,
the exact sixth derivative has the form
\[
 G_0^{(6)}(N)=\frac{H_6(r,\sigma)}{N^5L_N},
 \qquad \operatorname{den}H_6=(4+4r-3\sigma)^{12}.
\]
The numerator is a degree-thirteen polynomial with 82 terms.  An exact
coefficientwise majorant on $|r|,|\sigma|\le7/50$, together with
\eqref{eq:reducedden}, is less than $10^4$.  This gives the conservative
uniform bound $|G_0^{(6)}(N)|\le 20000/(|N|^5\log|N|)$.

The rational differentiation tower and its coefficientwise majorant are
machine-checked artifacts.  We do not duplicate that decisive symbolic
algebra here; Section~\ref{sec:mathematica} records its exact independent-CAS
reconstruction and frozen evidence.

\section{Six coefficient matches and the elementary map}\label{sec:model}

Use positive parameters
\[
 A=\frac{\alpha}{xe},\quad B=\frac{t+we}{x},\quad
 C=\frac{t}{x},\quad D=\frac{1+\delta e}{x},
 \qquad x=n^{-1},\ e=\cL_n^{-1}.
\]
The compact box is
\[
 \alpha\in[5/2,7/2],\quad t\in[7/4,9/4],\quad
 w\in[5,6],\quad\delta\in[1/4,5/12].
\]

For $f_k(U)=\log(U+k)-\log(U+k+1)$, the exact gamma half-shift is
$-f_k(n+1/2)$.  Pairing that term before estimation avoids a spurious
$\cL_n$ loss.  Repeated FTC gives, with $q=j+1$,
\[
 \Delta^jf_0(U)=(-1)^{j+1}j!
 \int_{[0,1]^q}(U+u_1+\cdots+u_q)^{-q}\,\dd u.
\]
Writing
\[
 \Phi_q(s,z)=\int_{[0,1]^q}
 (s+z\textstyle\sum u_i)^{-q}\,\dd u,
\]
one has, for $r=0,1,2$,
\[
 \partial_s^r\Phi_q(s,z)=(-1)^r(q)_r
 \int_{[0,1]^q}(s+z\textstyle\sum u_i)^{-q-r}\,\dd u.
\]
If $s\ge s_0>0$ and $z\ge0$, the unit cube has volume one and hence
\[
 |\partial_s^r\Phi_q(s,z)|\le(q)_r s_0^{-q-r}.
\]
The mean-value theorem adds the explicit first-order error
$(q)_r(q+r)qzs_0^{-q-r-1}$.

With $(a_0,a_1,a_2,a_3)=(1,1/2,1,1)$, the normalized elementary term is
\begin{align*}
H_{n,j}=\frac{a_j}{e}\{&e^q\Phi_q(\alpha,xe)
-\Phi_q(t+we,x)+\Phi_q(t,x)\\
&-\Phi_q(1+\delta e,x)+\Phi_q(1+x/2,x)\}.
\end{align*}
The $B-C$ divided difference tends to $qw/t^{q+1}$; the paired $D$--gamma
divided difference tends to $q\delta$; and the remote boundary contributes
$1/\alpha$ only for $q=1$.  Values and first parameter derivatives converge
uniformly, giving
\[
H^\infty=\left(
\frac1\alpha+\frac w{t^2}+\delta,
\frac w{t^3}+\delta,
\frac{3w}{t^4}+3\delta,
\frac{4w}{t^5}+4\delta\right)
\]
with error $O_K(e+x/e)$ in value and $O_K(e+x)$ in derivatives.
Combining \eqref{eq:tower} with the exact scales gives the xi contribution
$S^\infty=(-2,-1,-2,-2)$.

\begin{lemma}[Quotient adapter]\label{lem:quotient}
Let $(a_j)$ and $(b_j)$ be nonzero coefficients with compatible logarithm
branches.  Equality of the four consecutive quotient coordinates is equality
of the four second differences
\[
 (a_{k+2}-2a_{k+1}+a_k)=(b_{k+2}-2b_{k+1}+b_k),\quad0\le k\le3,
\]
in logarithmic coordinates.  Together with the two normalizations at indices
$0$ and $1$, these equations imply equality at all six indices $0,\ldots,5$.
\end{lemma}

\begin{proof}
For the difference $c_k=a_k-b_k$, the hypotheses say $c_0=c_1=0$ and
$c_{k+2}=2c_{k+1}-c_k$.  Induction gives $c_2=\cdots=c_5=0$.
\end{proof}

\section{The positive parameter branch}\label{sec:branch}

The limiting residual map $F=H^\infty+S^\infty$ has the positive zero
\[
 y_*=(\alpha,t,w,\delta)=(3,2,16/3,1/3).
\]
Exact elimination gives positive-orthant uniqueness.  Its Jacobian satisfies
\[
 \det DF(y_*)=-\frac1{144},\qquad
 \|DF(y_*)^{-1}\|_\infty=\frac{304}{3}.
\]
On a smaller rational box $K_0$ around $y_*$, choose
$P=DF(y_*)^{-1}$ so that $\|I-PDF\|_\infty\le1/4$.  The $C^1$ estimate in
Section~\ref{sec:model} makes $T_n(y)=y-PG_n(y)$ a self-map and a contraction
with constant at most $1/2$.  Thus there is a unique zero in $K_0$ and
\[
 y_n-y_*=O(\cL_n^{-1}).
\]
Uniqueness is asserted only in $K_0$, not globally.

For $0<e\le1/12$ and $x>0$, exact endpoint arithmetic gives
\[
 \frac{\alpha}{e}\ge30>\frac{11}{4}\ge t+we,\qquad
 t-1-\delta e\ge\frac{103}{144}>0.
\]
Consequently $A>B>C>D>0$ on the box.  The point $y_*$ has exact left/right
margins $(1/2,1/2)$, $(1/4,1/4)$, $(1/3,2/3)$, and $(1/12,1/12)$.

\section{Finite-free comparison and localization}\label{sec:finitefree}

The comparison polynomial is
\[
 p_F(y)={}_3F_2\!\left(
 \begin{matrix}-d,A,C\\B,D\end{matrix};\frac{D}{AC}y\right).
\]
For
\[
 Q_{U,V}^{(s)}(X)={}_2F_1\!\left(
 \begin{matrix}-d,U\\V\end{matrix};\frac{sX}{U}\right),
\]
direct multiplication of the terminating coefficients gives, in
constant-term-one normalization,
\[
 p_F=Q_{A,B}^{(1)}\boxtimes_d Q_{C,D}^{(D)}.
\]
Indeed, its $k$th coefficient is exactly
\[
 (-1)^k\binom dk
 \frac{(A)_k(C)_k}{(B)_k(D)_k}
 \left(\frac{D}{AC}\right)^k.
\]
Moreover
\[
 Q_{U,V}^{(s)}(X)=\frac{d!}{(V)_d}
 P_d^{(V-1,U-V-d)}\!\left(1-\frac{2sX}{U}\right).
\]
The branch inequalities $B,D>0$, $A\ge B+d$, and $C\ge D+d$ therefore put
both Jacobi parameters above $-1$ and make the two factors degree $d$ with
simple positive roots.  Reversal converts the coefficient identity to
the monic descending convention in the finite-free literature.  To pin the
convention, order positive roots as
$\lambda_1\ge\cdots\ge\lambda_d>0$ and, following MMP Definition~2.16, put
\[
 \operatorname{lmesh}(p)=
 \min_{1\le j<d}\frac{\lambda_j}{\lambda_{j+1}}\ge1.
\]
For $d\ge2$, the first Jacobi factor has logarithmic mesh strictly greater
than one.  MMP Proposition~2.7(iii) preserves nonnegative
real-rootedness, while Proposition~2.17 states in this convention that
\[
 \operatorname{lmesh}(p\boxtimes_dq)
 \ge \operatorname{lmesh}(p).
\]
Thus the convolution is simple.  Since $p_F(0)=1$, zero is not a root; degrees
zero and one are immediate.  The coefficient identity, rather than an
assumption about the final ${}_3F_2$, transports the result to $p_F$.  The
exact proposition text is pinned to the accessible arXiv v3; the paywalled
journal PDF has not been byte-compared.

For a Jacobi factor $q_{U,V}$, direct Gershgorin estimates on the transported
Jacobi matrix give the ratio-free statement
\[
 U\ge V+d,\quad V\ge32d
 \quad\Longrightarrow\quad
 \operatorname{roots}(q_{U,V})\subset
 [V-8\sqrt{Vd},V+8\sqrt{Vd}].
\]
Before deriving the product constant, the parameter branch and wedge give
the non-circular coarse box
\[
 n\le B\le3n,\qquad \frac n2\le D\le2n,
 \qquad A\ge8B,\qquad C\ge D+d,
 \qquad B,D\ge256d.
\]
Thus both Jacobi lower endpoints are positive and $B/D\le6$.  The
maximum-root product inequality is cited directly as MSS Theorem 1.6;
applying it also to reciprocal polynomials supplies the lower endpoint.  If
$u$ and $v$ denote the deviations of the two factor roots from $B$ and $1$,
respectively, then
\begin{align*}
 |(B+u)(1+v)-B|
 &\le8\sqrt{Bd}+8B\sqrt{d/D}
       +64\sqrt{Bd}\sqrt{d/D}\\
 &\le(12+8\sqrt6)\sqrt{Bd}.
\end{align*}
The resulting comparison roots satisfy
\begin{equation}\label{eq:local}
                    |y-B|\le C_{\rm loc}\sqrt{Bd},\qquad
 C_{\rm loc}=12+8\sqrt6.
\end{equation}
Critical points satisfy the same interval by interlacing.  The elementary
bound $\sqrt6<5/2$ gives $C_{\rm loc}<32$ and the exact rational sufficient
choice $K_0=256\cdot32^2=262144$ for the strengthened box.

This argument does not use Holland's assembled Lemma 7.3: its hypothesis on
$(C-D)/D$ fails on the present branch.  It also corrects a theorem-number
seam: the published MSS largest-root product inequality is Theorem 1.6.
In the typed formal boundary, the two factor positive-root and degree
certificates and the strict inequalities $0<u_-$ and $0<v_-$ are explicit
arguments.  Lean derives $u_-\ge B/2>0$ and $v_-\ge1/2>0$ from
$B,D\ge256d$ before invoking MSS.

\section{The direct critical-point recurrence}\label{sec:recurrence}

For the $m$th derivative, shift the hypergeometric parameters to
$(m-d,A+m,C+m;B+m,D+m)$.  The coefficient ratio proves directly
\[
 \vartheta(\vartheta+B+m-1)(\vartheta+D+m-1)g_m
 =\lambda y(\vartheta+m-d)(\vartheta+A+m)(\vartheta+C+m)g_m,
\]
Here
\[
 \vartheta=y\,\dd/\dd y,\qquad \lambda=D/(AC).
\]
Expanding the Euler operators gives, for
$T_k=y^kp_F^{(k)}(y)/p_F(y)$,
\[
 P_{3,m}T_{m+3}+P_{2,m}T_{m+2}
 +P_{1,m}T_{m+1}+P_{0,m}T_m=0.
\]
For a readable display of all four exact coefficients, set
\begin{align*}
 a&=1-y/A,&
 b_m&=B+m-y+(d-1-2m)y/A,\\
 c_m&=(d-m)(1+m/A)y,&
 \epsilon_p&=(C-D)/C,\\
 \beta_m&=A(2m+1-d)-d(2m+1)+3m^2+3m+1,&
 \Gamma_m&=m(m-d)(m+A).
\end{align*}
Then
\begin{align*}
 P_{3,m}&=a+\epsilon_p\frac yA=\frac{AC-Dy}{AC},\\
 P_{2,m}&=b_{m+1}+(D+m)a
   +\epsilon_p\frac yA(A-d+3+3m),\\
 P_{1,m}&=c_{m+1}+(D+m)b_m
   +\epsilon_p\frac yA\beta_m,\\
 P_{0,m}&=(D+m)c_m+\epsilon_p\frac yA\Gamma_m
   =\frac{Dy(A+m)(C+m)(d-m)}{AC}.
\end{align*}
The expanded rational forms are also checked by Lean; the shifted
hypergeometric ODE and these decompositions are checked by the exact
symbolic producer.

At a critical point, $T_0=1$ and $T_1=0$.  Localization
\eqref{eq:local}, the wedge, and positivity of the non-perturbative
$\epsilon_p=(C-D)/C$ term yield, uniformly for $0\le m\le d-2$,
\[
 P_{2,m}\ge n/8,\quad |P_{3,m}|\le2,\quad
 |P_{1,m}|\ll n\sqrt{Bd}+nd,\quad 0\le P_{0,m}\ll n^2d.
\]
A maximum argument in the recurrence then gives the critical-point
derivative radius
\begin{equation}\label{eq:radius}
 \max_{1\le k\le d}
 \left|\frac{y^kp_F^{(k)}(y)}{p_F(y)}\right|^{1/k}
 \le K_r\sqrt{Bd}
\end{equation}
with an absolute $K_r$; separately $T_0=1$.  This is the decisive replacement for a perturbative
recurrence that is unavailable when $\epsilon_p\to1/2$.

\section{Sixth-order residual}\label{sec:residual}

Let $E_F(z)$ be the logarithm of the exact xi/model coefficient ratio along
the positive branch.  Six differentiations give
\begin{align*}
E_F^{(6)}(z)={}&h^{(6)}(n+z)+\psi^{(5)}(B+z)
-\psi^{(5)}(n+1/2+z)-\psi^{(5)}(A+z)\\
&+\psi^{(5)}(D+z)-\psi^{(5)}(C+z).
\end{align*}
Pair the $B-C$ and $D-(n+1/2)$ terms.  On the thickened interpolation domain
\[
 \Omega=\{z:\operatorname{dist}(z,[0,d])\le2\rho\},\qquad
 \rho=K_r\sqrt{Bd},
\]
the wedge gives $d+2\rho\le\eta n$ for a fixed small $\eta$.  Hence
$n+\Omega$ remains in the nonvanishing coefficient sector, while every
straight polygamma segment stays in $\Re z\ge cn$.  The absolutely convergent
polygamma series and Section~\ref{sec:deriv} imply
\begin{equation}\label{eq:E6}
 \sup_{z\in\Omega}|E_F^{(6)}(z)|
 \le \frac{C}{n^5\log(n+2)}.
\end{equation}

The complex Hermite--Genocchi formula is obtained by iterating the complex
line-segment FTC over the standard six-simplex.  Its mass is $1/6!=1/720$.
Since $E_F(0)=\cdots=E_F(5)=0$, Newton interpolation and
\eqref{eq:E6} give, for the needed $z$,
\begin{equation}\label{eq:HG}
 |E_F(z)|\le \frac{C}{720n^5\log(n+2)}
       \prod_{j=0}^{5}|z-j|.
\end{equation}
The wedge implies $\rho\ge d$ after increasing its absolute constant.  For
$z\in\Omega$ and $0\le j\le5$, projection to $[0,d]$ then gives
$|z-j|\le3\rho$.  Since $B\asymp n$,
\[
 \rho^6=K_r^6(Bd)^3\ll n^3d^3,
\]
and \eqref{eq:HG} yields the headline exponent calculation
\begin{equation}\label{eq:exponent}
 \sup_{z\in\Omega}|E_F(z)|
 \ll\frac{\rho^6}{n^5\log(n+2)}
 \ll\frac{d^3}{n^2\log(n+2)}\le\frac{C'}K.
\end{equation}
Choose the absolute wedge constant so that the last quantity is small enough
that, for $c_F(z)=e^{E_F(z)}$,
$\sup_\Omega|c_F-1|\le1$.  Exponentiation is legal on the same chosen
logarithm branch.  At the integer nodes the exact coefficient ratios are
positive real numbers, and the six matches give
$c_F(0)=\cdots=c_F(5)=1$.

\section{Sign transfer and proof of the main theorem}\label{sec:assembly}

\begin{lemma}[Finite multiplier stability]\label{lem:stability}
Let $d\ge5$, and let $p(y)=\sum_{j=0}^dp_jy^j$ be real with $d$ simple
positive zeros and $p(0)\ne0$.  Put
$\Omega_r=\{z:\operatorname{dist}(z,[0,d])\le2r\}$, and let $c$ be
holomorphic on a neighborhood of $\Omega_r$.  Suppose that
$c(0),\ldots,c(d)$ are real,
$c(0)=\cdots=c(4)=1$, $c(d)>0$,
$\sup_{\Omega_r}|c-1|\le\epsilon<16$, and, at every critical point $y$ of
$p$,
\[
 \left|\frac{y^kp^{(k)}(y)}{p(y)}\right|\le r^k,
 \qquad0\le k\le d.
\]
Then $P(y)=\sum_{j=0}^dp_jc(j)y^j$ has $d$ distinct positive real zeros.
\end{lemma}

\begin{proof}
Newton interpolation and termwise differentiation give
\[
 \frac{P(y)}{p(y)}
 =\sum_{k=0}^d\frac{\Delta^kc(0)}{k!}
   \frac{y^kp^{(k)}(y)}{p(y)}.
\]
Cauchy's estimate on the radius-$2r$ discs in $\Omega_r$ gives
$|\Delta^kc(0)|/k!\le\epsilon/(2r)^k$.  The five exact matches make the
terms $1\le k\le4$ vanish, so at every critical point
\[
 \left|\frac{P(y)}{p(y)}-1\right|
 \le\epsilon\sum_{k=5}^d2^{-k}<\epsilon/16<1.
\]
Thus $P$ and $p$ have the same signs at all critical points, at zero, and
for large positive $y$.  Interlacing and the intermediate-value theorem put
one simple sign-changing zero of $P$ in each of the $d$ intervening
intervals.
\end{proof}

\begin{proof}[Proof of Theorem~\ref{thm:main}]
The comparison polynomial has $d$ simple positive roots.  Between them lie
$d-1$ critical points, and the two exterior sign tests complete the
alternating pattern.  For $d\le5$, the six coefficient matches already cover
the whole polynomial.  For $d\ge6$, apply Lemma~\ref{lem:stability} with
$p=p_F$, $c=c_F$, and $r=\rho$.  Equations \eqref{eq:radius},
\eqref{eq:exponent}, and the positivity of the integer coefficient ratios
verify every hypothesis with $\epsilon\le1<16$.  The lemma therefore
preserves every comparison sign, so the
intermediate-value theorem supplies $d$ distinct positive roots of the
transformed actual polynomial.

The positive parameter branch has a positive scale and nonzero
normalization.  Undoing $y=-SX$, with $S>0$, carries those roots to $d$
distinct negative roots of $\J^{d,n}$.  This proves
Theorem~\ref{thm:main}.  If $N_0$ is the eventual analytic cutoff, set
$K_{\rm finite}=N_0^2(N_0+2)+1$.  For $n<N_0$ and $d\ge1$,
\[
 n^2\log(n+2)\le n^2(n+1)<K_{\rm finite}\le K_{\rm finite}d^3,
\]
so the wedge antecedent is empty.  Lean checks this arithmetic in the theorem
whose source name is
\texttt{not\_twoThirdsWedge\_finiteCutoffAbsorption}; no low-index
hyperbolicity is assumed.  A second kernel argument proves that the Jensen
polynomial has degree exactly $d$ and cannot possess $d+1$ distinct roots.
The declaration
\texttt{riemannXiJensen\_twoThirds\_global\_headline\_exactly} combines the
analytic and pre-cutoff cases and returns exactly $d$ distinct negative roots.
\end{proof}

\section{Formal verification and trust boundary}\label{sec:formal}

The Lean project contains no axiom asserting the xi asymptotic or the headline
result.  It checks the complete concrete T1--T5 chain: the xi/theta/Mellin
identity, quantitative saddle branch, legal leading contour, infinite
higher-mode suppression, Gamma/Stirling assembly, sectorial coefficient
theorem, and Cauchy transport through order six.  It also checks the rational
leading system and Jacobian, quotient-to-six-coefficient adapter, cube
calculus, exact xi branch and six matches, direct recurrence, multiplier sign
transfer, concrete Jensen definition, and positive-to-negative scaling.

The complex Hermite--Genocchi modules prove complex line-segment FTC, six
repeated FTC steps, the exact Newton identity, convex-hull localization, the
six-simplex mass $1/720$, and the local remainder bound.  The approved
Phase-30 endpoint constructs the concrete xi multiplier, proves the first six
node values and its strict unit bound, instantiates the complete
Rolle-point interval certificate, and connects sign transfer to the actual
transformed Jensen polynomial.  Its MMP record is now attached to the two
concrete Jacobi factors, including their positive-root and degree hypotheses;
Lean uses the coefficient identity above to transport the MMP conclusion to
the actual xi comparison.  Classical Jacobi root/matrix facts, general MMP
positivity/log-mesh preservation, and the MSS product-root inequality remain
narrow typed literature inputs.  These limitations are proof-surface
disclosures, not hidden axioms.

The MSS record also carries the two factor-domain certificates and requires
strictly positive lower endpoints.  A negative regression reproduces the
formerly unguarded call and must fail at the missing positivity proof.  The
current terminal theorem performs the finite-cutoff split in one declaration
and upgrades the root statement from at least $d$ to exactly $d$ using the
kernel-checked degree calculation.

The literal T5 statements have been isolated in a locally verified Phase-29
Palomar submission candidate.  Official Comparator and NanoDa replay remains
pending, so no official Palomar result is claimed.

The verifier scans for \texttt{sorry}, \texttt{admit}, custom axioms, and
unsafe escape hatches, builds the headline target with the pinned toolchain,
and audits selected theorem axioms.  Numerical contour and Mellin checks are
regression evidence only.

\section{Independent-CAS Mathematica reconstruction}\label{sec:mathematica}

A user-executed Mathematica 15.0.1 reconstruction returned exact matches for
M1--M4.  The user entered and executed cells supplied interactively by Codex;
the notebook imported no repository JSON, Python, SymPy, Lean-generated
expression, or frozen numerator.  Thus the run breaks the
SymPy-versus-Mathematica CAS common mode, while remaining AI-assisted evidence
rather than human or peer review.

The exact terminal results were: the implicit saddle derivative identity and
post-chain constants $(2,-2,4,-12,48)$; the shifted ${}_3F_2$ coefficient
recurrence, four zero coefficient differences, and complete derivative-order
tests at degrees $5,8,11$; the positive solution
$(3,2,16/3,1/3)$, determinant $-1/144$, and inverse norm $304/3$; and the
82-term degree-13 $H_6$ numerator with coefficientwise majorant
\[
 \frac{6422139805764931584036533551104}
 {702576099728137594188684005}<10^4.
\]
All mathematical outputs are exact and contain no machine-real values.  The
notebook, plain-text result ledger, visual PDF, and user checksum ledger are
frozen in the Phase-24 evidence directory.  The canonical result-ledger
SHA-256 is
\texttt{1faea5fc\allowbreak b35b504b\allowbreak 5d0aad93\allowbreak
91999a99\allowbreak e530373d\allowbreak ce36addb\allowbreak
496eabb1\allowbreak 33fb04cc}.

\section{Availability and review statement}

The Version 1.0 public-release set is designed to contain this source and rendered PDF,
all proof notes needed to expand the compressed arguments, the Lean source
and lockfiles, exact symbolic and interval certificates, verifier scripts,
serial logs, primary-source hashes, and AI-review disclosures.  Public
posting has been authorized, but the absence of human mathematical or peer
review remains a material limitation and is disclosed on every circulation
surface.  The public version should be cited by its version-specific DOI and
checksum once registered; corrections should appear as new versions rather
than silent file replacement.

\begin{thebibliography}{9}
\bibitem{GORZ}
M. Griffin, K. Ono, L. Rolen, and D. Zagier,
\emph{Jensen polynomials for the Riemann zeta function and other sequences},
Proceedings of the National Academy of Sciences 116 (2019), 11103--11110.

\bibitem{GORTTW}
M. Griffin, K. Ono, L. Rolen, J. Thorner, Z. Tripp, and I. Wagner,
\emph{Jensen polynomials for the Riemann xi-function},
arXiv:1910.01227v3.

\bibitem{Holland}
J. Holland,
\emph{A new hyperbolicity wedge and a joint semicircle limit for Jensen
polynomials of Riemann's xi-function}, arXiv:2608.08682v1.

\bibitem{MMP}
A. Martinez-Finkelshtein, Rafael Morales, and D. Perales,
\emph{Real roots of hypergeometric polynomials via finite free convolution},
International Mathematics Research Notices 2024 (2024), 11642--11687;
arXiv:2309.10970v3.

\bibitem{MSS}
A. Marcus, D. Spielman, and N. Srivastava,
\emph{Finite free convolutions of polynomials},
Probability Theory and Related Fields 182 (2022), 807--848.

\bibitem{Szego}
G. Szego, \emph{Orthogonal Polynomials}, fourth edition,
American Mathematical Society Colloquium Publications 23, 1975.
\end{thebibliography}

\end{document}
